% !TeX program = xelatex
% ==========================================================================
% Пропоузал семестрового проєкту — курс «Штучний інтелект», УКУ.
% Це ЄДИНИЙ файл, який вам потрібно редагувати. Уся верстка й брендинг
% живуть у ucuproposal.cls — не змінюйте шрифти, поля, кеглі й кольори.
%
% Компіляція: XeLaTeX + Biber (Overleaf робить це автоматично).
% Локально:   latexmk (див. README.md).
% ==========================================================================
\documentclass[ukrainian]{ucuproposal}
% Опції класу (можна комбінувати через кому):
%   ukrainian | english  — мова документа й логотипу (за замовчуванням ukrainian)
%   slogan                — додає гасло УКУ в колонтитул титулки (за замовчуванням вимкнено)

% --------------------------------------------------------------------------
% Метадані титульної сторінки
% --------------------------------------------------------------------------
\projecttitle{Аналіз алгоритмів unlearning в LLM для задач bias mitigation}
\projecttagline{Частина проєкту FairForget --- інструментарію для перевірки моделі на упередженість до окремих спільнот.}

\teammember{Гаєвська Юстина}
\teammember{Ружило Павло}
\teammember{Копач Тарас}
% \teammember{Прізвище Ім'я} % третій учасник (опційно)
% \teammember[https://linkedin.com/in/...]{Прізвище Ім'я} % опційне клікабельне посилання на профіль (LinkedIn тощо)

\mentor{Маковська Вікторія} % якщо ментора немає — видаліть текст і залиште \mentor{}, рядок не друкується

\track{research} % research | engineering — не друкується на титулці, лише для внутрішньої категоризації

% \repo{https://github.com/team/project} % опційно; якщо репозиторію ще немає — залиште порожнім, рядок не друкується

\proposaldate{2026\mbox{-}09\mbox{-}20}

\begin{document}
\maketitlepage{}

\section{Проблема та мотивація}
% Що саме ви розв'язуєте і чому це важливо й цікаво в контексті ШІ?
% Орієнтовно 1 абзац.
З поширенням використання інструментів штучного інтелекту у повсякденному житті, виразнішою стає проблема небажаної поведінки мовних моделей: відтворення шкідливої інформації, порушення приватності, генерація небезпечного контенту або відтворення стереотипних та дискримінаційних упереджень. Звідси породжується потреба "відівчити" модель від даних, що викликають цю поведінку. Найбільш очевидним способом є повторне навчання моделі на даних, обмежених від небажаних. Проте для великих мовних моделей (LLM) як ChatGPT чи Claude це може вимагати значних фінансових витрат та часу. На цьому фоні розвивається напрям Machine Unlearning: дослідження алгоритмів, що обмежують небажану поведінку моделей без перенавчання. Сфера нашої діяльності базується навколо зменшення саме стереоптипних упереджень мовної моделі (Bias mitigation), використовуючи алгоритми unlearning.


\section{Огляд джерел та наявних рішень}
% Що вже зроблено у світі за цією темою? Спирайтесь на references.bib,
% цитуйте через \cite{}. Орієнтовно 3-5 джерел.
Representation Misdirection for Unlearning (RMU) \cite{pmlr-v235-li24bc} — це метод Machine Unlearning, який змінює внутрішні представлення моделі для даних із forget set, спрямовуючи їх до випадкового цільового представлення. Одночасно модель оптимізується так, щоб зберегти представлення даних із retain set. Таким чином, RMU намагається усунути вплив небажаної інформації без повного перенавчання моделі. 
- 


*Методи з туторіалу, які це роблять, референс пейпери*
- Vanilla \& adaptive RMU~\cite{pmlr-v235-li24bc}\cite{dang2025effects}

\section{Дані}
% Які набори даних, звідки, ліцензія. Якщо збираєте самі — як саме.
% Якщо дані не потрібні — поясніть чому.
WarBias \cite{makovska_etal_2026_warbias} -- синтетична збірка складена зі стереоптипів, контрстереотипів та нейтральних висловів в контексті війни. Загальна кількість записів становить 5,238 українською та англійською. Ліцензія: CC BY 4.0.

BBQ-UK: Ukrainian Translation\cite{makovska_etal_2026_bbq_uk} -- український переклад бенчмарки Bias Benchmark for Question Answering (BBQ).\cite{parrish-etal-2022-bbq}
Даний набір має в собі неоднозначні питання, які можуть мати стереотипічні відповідді. Загальна кількість записів: 57,006. Варто зазначити, що хоч і BBQ-UK має в собі суто українські переклади, у ньому зберігаєтьяс переважний американський побутовий контекст. Ліцензія: CC BY 4.0.
\section{Підхід і методи}
% Які алгоритми/архітектури плануєте використати, чи спираєтесь на
% існуючі реалізації та як плануєте їх покращити.
Задачу machine unlearning можна формально описати наступним чином: \\
Нехай $M(X)$ — мовна модель $M$, навчена на множині даних $X$, а $S \subseteq X$ — множина даних, пов'язаних із небажаною поведінкою моделі.

Задача Machine Unlearning полягає у знаходженні алгоритму ($U$), який за моделлю ($M(X)$) та множиною даних ($S$) будує модель, поведінка якої наближається до поведінки моделі, отриманої шляхом навчання виключно на ($X \setminus S$):

\[
  U(M(X), S) \sim M(X \setminus S).
\]
Щоб досягти цього, плануємо використовувати принцип \textbf{R}epresentation \textbf{M}isdirection In \textbf{U}nlearning та можливі модифікації (Adaptive RMU).
Алгоритм досягає виконання задачі mitigation/unlearning за допомогою зниження точності на датасеті з стереотипними питаннями, які треба "забути", водночас зберігаючи хорошу точність на датасеті, що тестує загальні можливості моделі (т.зв \textit{retain dataset}, до прикладу, LLMU). Ключовим принципом роботи алгоритму є використання двочастинної функції витрат:
\[
L = L_{\mathrm{forget}} + L_{\mathrm{retain}}
\]
\section{Оцінювання результатів}
% Метрики, baseline, спосіб порівняння, якісні очікування.
Опишіть метрики та baseline, з яким порівнюватиметесь.
WMDP i LLMU MT-bench
\section{План і ризики}
% Етапи до midterm і до фінального захисту, ключові ризики та план Б.
Опишіть план роботи та основні ризики.

% Приклад врізки з ключовими цифрами (опційно, можна видалити):
% \begin{keyfacts}
% \textbf{Ключові цифри:} 3 датасети \textperiodcentered\ 2 baseline-моделі \textperiodcentered\ 4 тижні до midterm
% \end{keyfacts}

\section{Contribution statement}
% Хто за що відповідає. Відсотки не обов'язкові — якщо хочете їх вказати,
% додайте в квадратних дужках: \contribution[50\%]{Прізвище Ім'я}{...}
\begin{contributions}
** TBD? ми цього не знаємо або можемо запитати в чат
  \contribution{Прізвище Ім'я}{дані, препроцесинг, EDA}
  \contribution{Прізвище Ім'я}{моделювання, евалюація, документація}
\end{contributions}

% Примітка: декларація використання AI-інструментів (AI usage) до цього
% пропоузалу не входить — вона живе в README.md вашого репозиторію.

\ucuprintbibliography{}

\end{document}
